\documentclass[11pt,a4paper]{article}
\usepackage[a4paper,margin=2.4cm]{geometry}
\usepackage{fontspec}
\usepackage{amsmath,amssymb,amsthm}
\usepackage{unicode-math}
\setmainfont{Liberation Serif}
\setsansfont{Liberation Sans}
\setmonofont{DejaVu Sans Mono}[Scale=0.85]
\setmathfont{Latin Modern Math}
\usepackage{polyglossia}
\setdefaultlanguage{russian}
\setotherlanguage{english}
\newfontfamily\cyrillicfont{Liberation Serif}
\newfontfamily\cyrillicfontsf{Liberation Sans}
\newfontfamily\cyrillicfonttt{DejaVu Sans Mono}[Scale=0.85]
\usepackage{booktabs}
\usepackage{enumitem}
\setlist{nosep,leftmargin=2em}
\usepackage[hidelinks]{hyperref}
\usepackage{microtype}

\theoremstyle{plain}
\newtheorem{theorem}{Теорема}
\newtheorem{lemma}[theorem]{Лемма}
\newtheorem{proposition}[theorem]{Предложение}
\newtheorem{corollary}[theorem]{Следствие}
\theoremstyle{definition}
\newtheorem{definition}{Определение}
\newtheorem{conjecture}{Гипотеза}
\newtheorem{problem}{Проблема}
\theoremstyle{remark}
\newtheorem*{remark}{Замечание}
\newtheorem*{example}{Пример}

\newcommand{\Z}{\mathbb{Z}}
\newcommand{\R}{\mathbb{R}}
\newcommand{\lo}{\lambda}
\DeclareMathOperator{\vol}{vol}

\title{\textbf{Относительная сетка:\\ действия над клеточными фигурами,\\ потеря информации и число действий}}
\author{М.~Зафаров\thanks{Автор системы, определений и гипотез. Доказательства и текст оформлены с~помощью ИИ-помощника Claude (Anthropic). Текст распространяется по лицензии CC~BY~4.0; программы и~файлы Lean — по лицензии MIT. Полный текст, программы проверки и формальная запись: \url{https://github.com/abema2416-ops/Relative-Grid}.}}
\date{Препринт, 27 ноября 2026 г.}

\begin{document}
\maketitle

\begin{abstract}
Рассматривается система действий над конечными множествами клеток решётки $\Z^n$: растяжение и сжатие части объекта вдоль одной оси в целое число раз ($\alpha$), удаление клеток ($\sigma$), склеивание и расклеивание касающихся клеток ($\beta_\pm$). Каждому действию сопоставляется потеря информации — двоичный логарифм числа состояний зоны, переходящих в один и тот же результат, — а каждому объекту число действий $\ell(O)$, необходимых для его построения из одной клетки. Доказано: каждая исчезнувшая клетка стоит не меньше одного бита, а для сжатия ряда найдено точное число прообразов; клеточное растяжение совпадает с аффинным растяжением пространства без округлений; на плоскости растяжения не меняют число дыр, а одна склейка меняет его не больше чем на единицу. Для «рельефов» вычисление числа действий содержит NP-трудную задачу о нормальном базисе множеств; получены оценки и точные частные случаи. Сформулирована гипотеза («правило отрезков»), сводящая перебор всех зон к перебору отрезков линий объекта; она проверена полным перебором на $4\,835\,965$ фигурах в размерностях 2 и~3. Ядро системы записано и проверено в Lean~4.
\end{abstract}

\section{Введение}

Фигуры, нарезанные сеткой на клетки, — обычный объект дискретной геометрии, цифровой топологии и обработки изображений. Над ними обычно рассматриваются морфологические операции (дилатация, эрозия) \cite{serra}, перемещения отдельных клеток \cite{akitaya} или самосборка \cite{soloveichik}. В настоящей работе предлагается иной набор элементарных действий: часть объекта растягивается или сжимается вдоль одной оси в целое число раз, как гармошка, от своего крайнего слоя. Такое действие точно соответствует аффинному растяжению пространства (теорема~\ref{thm:affine}), но параметры (зона, ось, направление, коэффициент) задаются в клетках.

Вводятся две меры.
\emph{Потеря} действия — число бит, которые нельзя восстановить по результату и параметрам; это мера логической необратимости в смысле Ландауэра \cite{landauer}, применённая к действиям над формами.
\emph{Число действий} $\ell(O)$ — длина кратчайшей цепочки, строящей объект $O$ из одной клетки; это аналог описательной сложности формы \cite{soloveichik} и индекса сборки \cite{assembly}, но для геометрических действий.
Отдельно от касаний вводятся склейки: две клетки могут касаться, не будучи соединены. Это позволяет считать дыры как функцию выбранных склеек (теорема~\ref{thm:holes}).

Основные результаты: нижняя оценка потери и точная формула для сжатия ряда (§\ref{sec:loss}); мост к непрерывной геометрии (§\ref{sec:affine}); склейки и дыры на плоскости (§\ref{sec:holes}); оценки числа действий и NP-трудность его счёта для рельефов (§\ref{sec:count}); гипотезы, подкреплённые полным перебором, и открытые проблемы (§\ref{sec:conj}). В §\ref{sec:formal} описаны формальная проверка в Lean~4 и программы.

\section{Определения}\label{sec:def}

Обозначим через $e_1,\dots,e_n$ единичные векторы осей, $|X|$ — число элементов конечного множества $X$, $\|v\|_1=\sum|v_i|$, $\|v\|_\infty=\max|v_i|$.

\begin{definition}[клетка, объект]
\emph{Клетка} — точка $c\in\Z^n$; она изображает куб $Q(c)=c+[0,1]^n$. \emph{Объект} — пара $(O,S)$, где $O\subset\Z^n$ — конечное непустое множество клеток, а $S$ — множество \emph{склеек}: неупорядоченных пар клеток из $O$, которые касаются ($\|c-c'\|_\infty=1$). Касание без склейки называется \emph{трещиной}. Объекты, отличающиеся сдвигом на вектор из $\Z^n$, считаются одной формой. \emph{Куском} называется компонента связности графа склеек.
\end{definition}

\begin{definition}[действие $\alpha$]
Параметры: конечная зона $Z\subset\Z^n$, ось $i$, направление $s\in\{\pm1\}$, коэффициент $f\in\{2,3,\dots\}$ и вид («растянуть» или «сжать»). Пусть $G=O\cap Z\neq\varnothing$ — группа. Положим $a=\min\{s c_i: c\in G\}$ (закреплённая сторона), $t(c)=s c_i-a\ge0$, и для целого $p$ обозначим через $c[p]$ клетку, у которой $i$-я координата равна $sp$, а остальные — как у $c$. Тогда
\[
\text{растяжение: } \alpha(O)=(O\setminus G)\cup\{c[a+f\,t(c)+j]: c\in G,\ 0\le j<f\},
\]
\[
\text{сжатие: } \alpha(O)=(O\setminus G)\cup\{c[a+\lfloor t(c)/f\rfloor]: c\in G\}.
\]
Действие допустимо, если новые клетки не пересекаются с $O\setminus G$ (правило П3). Образ клетки при растяжении — $f$ клеток подряд, склеенных между собой; склейка $\{c,c'\}$ переходит в склейки всех касающихся по стороне пар их образов. Растяжение или сжатие, после которого склеенные клетки перестали бы касаться, недопустимо (правило П5): склейку можно убрать только действием $\beta_-$ или вместе с клеткой.
\end{definition}

\begin{definition}[$\sigma$, $\beta_\pm$]
$\sigma_Z(O)=O\setminus Z$. Действие $\beta_+$ добавляет к $S$ непустое множество пар касающихся, но не склеенных клеток; $\beta_-$ удаляет непустое множество склеек. Клетки при этом не меняются.
\end{definition}

\begin{definition}[потеря, число действий]
Действие $\varphi$ с фиксированными параметрами применимо к любому содержимому своей зоны. \emph{Потерей} $\lo(\varphi,O)$ называется $\log_2 N$, где $N$ — число объектов $O'$, совпадающих с $O$ вне зоны, для которых $\varphi$ допустимо и $\varphi(O')=\varphi(O)$. Потеря цепочки — сумма потерь шагов. \emph{Число действий} $\ell(O)$ — наименьшая длина цепочки действий, переводящей одну клетку в $O$ (с точностью до сдвига).
\end{definition}

Растяжение, $\beta_+$ и $\beta_-$ потерь не дают: по результату и параметрам исходное состояние восстанавливается однозначно (для растяжения — потому что образы разных клеток не пересекаются).

\begin{definition}[тонкие объекты]
$k$-клетка ($0\le k\le n$) — множество $c+\sum_{j\in I}[0,1]e_j$, $|I|=k$. Объект разряда $k<n$ — конечное множество $k$-клеток. Действие $\alpha$ над таким объектом есть действие $\alpha$ пространства, применённое к каждой его клетке как к множеству точек: отображение $y\mapsto a+f(y-a)$ по оси $i$. Клетка, вытянутая вдоль оси $i$, переходит в $f$ клеток разряда $k$; клетка, «тонкая» вдоль оси $i$, — в одну клетку, сдвинутую в точку $a+f(c_i-a)$. Объект разряда $k>n$ в пространстве разряда $n$ изображается только частично.
\end{definition}

Тонкие объекты используются ниже лишь в примерах; все теоремы относятся к объектам полного разряда $k=n$.

\section{Потеря информации}\label{sec:loss}

\begin{lemma}\label{lem:bit}
Если действие уменьшает число клеток на $d>0$, то его потеря не меньше $d$ бит.
\end{lemma}
\begin{proof}
Растяжение клеток не убирает. Для $\sigma_Z$ результат не зависит от содержимого зоны, а содержимых не меньше $2^d$. Пусть действие — сжатие. Клетки группы в одном ряду вдоль оси $i$ делятся на блоки по значению $\lfloor t(c)/f\rfloor$; клетки блока сливаются в одну. Если в непустом блоке $B$ лежит $m_B$ клеток, то $d=\sum_B(m_B-1)$. Отметим клетку $e\in G$, на которой достигается $a$. Оставим в каждом непустом блоке любой непустой поднабор его клеток, причём $e$ — обязательно. Закреплённая сторона и номера $t(c)$ оставшихся клеток не меняются, каждый блок остаётся непустым, поэтому результат тот же и действие допустимо. Вариантов не меньше $2^{m_{B_0}-1}\prod_{B\ne B_0}(2^{m_B}-1)\ge 2^{d}$.
\end{proof}

\begin{theorem}\label{thm:loss}
При любом превращении объекта $A$ в объект $B$ суммарная потеря не меньше $|A|-|B|$ бит. Если $A$ связен по граням и $B=A\setminus H\neq\varnothing$, $|H|=h$, то наименьшая потеря превращения $A\to B$ равна ровно $h$ бит, а превращения $B\to A$ — нулю.
\end{theorem}
\begin{proof}
Число клеток меняется по шагам на $-d_1,-d_2,\dots$, и $\sum d_j=|A|-|B|$. По лемме~\ref{lem:bit} потеря шага не меньше $\max(d_j,0)$, сумма не меньше $|A|-|B|$. Значение $h$ достигается действием $\sigma_H$ ($2^h$ прообразов). Обратно, каждая компонента связности $H$ по граням граничит по грани с $B$, поэтому клетки $H$ можно упорядочить так, что каждая граничит с уже построенной; каждую возвращаем растяжением соседней клетки вдвое в её сторону, без потерь.
\end{proof}

Итак, убрать $h$ клеток стоит $h$ бит, вернуть — ничего. Это асимметрия, которой нет у числа действий: оба превращения требуют сравнимой работы.

\begin{theorem}[точная потеря сжатия ряда]\label{thm:exact}
Пусть зона — отрезок ряда с позициями $0,\dots,L-1$, $a$ — наименьшая занятая позиция, и ряд сжимается в $f$ раз. Разобьём позиции $a,\dots,L-1$ на блоки по $\lfloor (p-a)/f\rfloor$; $z_B$ — число позиций зоны в блоке $B$, $B_0$ — блок якоря. Тогда число содержимых зоны с тем же результатом равно
\[
N=2^{\,z_{B_0}-1}\prod_{B\ne B_0,\ B\text{ непуст}}\bigl(2^{z_B}-1\bigr).
\]
\end{theorem}
\begin{proof}
Результат сжатия определяется якорем $a$ и множеством непустых блоков. Содержимое $X$ даёт тот же результат ровно тогда, когда его самая левая клетка — $a$ (иначе сдвинется первая клетка результата или появится клетка левее $a$) и непустыми являются те же блоки. В блоке якоря клетка $a$ обязательна, остальные $z_{B_0}-1$ позиций свободны; в другом непустом блоке подходит любой непустой поднабор; пустые блоки пусты. Варианты независимы.
\end{proof}

\begin{example}
Для зоны длины 20, $f=2$, с четырьмя исчезнувшими клетками лемма даёт $N\ge16$, а теорема — $N=4374$. Формула совпала с полным перебором на $2872$ случайных рядах длины до 13 при $f=2,3,4$.
\end{example}

\section{Связь с аффинными растяжениями}\label{sec:affine}

Обозначим через $F(O)=\bigcup_{c\in O}Q(c)$ фигуру, изображаемую объектом, и пусть $y=sx_i$. Действие $\alpha$ задаёт отображение $T\colon y\mapsto a+f(y-a)$, остальные координаты неподвижны.

\begin{theorem}\label{thm:affine}
\begin{enumerate}[label=\upshape(\alph*)]
\item Для растяжения $F(\alpha(O))=F(O\setminus G)\cup T(F(G))$ — точно, без округления.
\item Для сжатия $\alpha(O)=(O\setminus G)\cup\{c: \vol(Q(c)\cap T^{-1}(F(G)))>0\}$, т.\,е. клеточное сжатие есть аффинное сжатие с последующей записью «клетка занята, если задета».
\item Пусть фигуры $F$ и $K\subset F$ ограничены, на уровнях сетки $N_1\mid N_2\mid\cdots$ записи $O_N$ и группы $G_N$ сходятся к ним по объёму симметрической разности, а якоря $a_N/N\to A$. Тогда растяжения $\alpha(O_N)$ сходятся по объёму к $(F\setminus K)\cup T_A(K)$; то же для любой конечной цепочки растяжений.
\end{enumerate}
\end{theorem}
\begin{proof}
(a) Клетка с позицией $p=a+t$ переходит в клетки с позициями $a+ft,\dots,a+ft+f-1$, которые вместе занимают по $y$ отрезок $[a+ft,\,a+ft+f)=T([p,p+1))$, а по остальным осям — то же, что исходная клетка. (b) $T^{-1}$ переводит куб клетки в слой толщины $1/f$ внутри клетки с позицией $a+\lfloor t/f\rfloor$, а слой имеет положительный объём. (c) По (a) и неравенству $\vol((X\cup Y)\triangle(X'\cup Y'))\le\vol(X\triangle X')+\vol(Y\triangle Y')$; $T$ умножает объёмы на $f$, а смещение якоря на $(f-1)(A-a_N/N)\to0$ даёт малый сдвиг фигуры конечного объёма.
\end{proof}

Таким образом, само действие $\alpha$ — классическое аффинное растяжение; новыми в системе являются способ задания (зона и закреплённая сторона в клетках), учёт потерь и счёт числа действий. Предел цепочек со сжатиями и $\sigma$ остаётся открытым (проблема~\ref{pr:limit}).

\section{Склейки и дыры на плоскости}\label{sec:holes}

Пусть $n=2$ и склейки только по сторонам. Построим двумерный клеточный комплекс $K(O,S)$: возьмём по квадрату $Q(c)$ на каждую клетку и отождествим общие стороны склеенных пар. Числом дыр назовём $d=b_1(K)$; через $q$ обозначим число точек решётки, вокруг которых стоит целиком склеенный квадрат $2\times2$.

\begin{theorem}\label{thm:holes}
\begin{enumerate}[label=\upshape(\alph*)]
\item $d=(\text{число кусков})-|O|+|S|-q$.
\item Допустимое растяжение не меняет $d$.
\item Добавление одной склейки меняет $d$ на $+1$, $0$ или $-1$.
\item Чтобы построить объект из одной клетки растяжениями и склейками, нужно склеить не меньше $d$ пар; при $q=0$ ровно $d$ пар достаточно.
\end{enumerate}
\end{theorem}
\begin{proof}
(a) Граней $F=|O|$. Каждая сторона участвует не больше чем в одной склейке, поэтому рёбер $E=4|O|-|S|$. Над точкой $p$ решётки рассмотрим граф $\Gamma(p)$: вершины — квадраты с углом $p$ ($k(p)$ штук), рёбра — склеенные стороны с концом $p$ ($s(p)$ штук). Это подграф цикла длины 4, поэтому число его компонент равно $k(p)-s(p)+[\Gamma(p)\text{ — весь цикл}]$; вершин $K$ над $p$ столько же. Суммируя ($\sum k=4|O|$, $\sum s=2|S|$), получаем $V=4|O|-2|S|+q$ и $\chi(K)=|O|-|S|+q$. Окрестность каждой точки $K$ — диск или полудиск (веер из $j\le4$ секторов по $90^\circ$ либо полный круг), значит $K$ — компактная поверхность, у каждой компоненты которой край непуст; поэтому $b_2=0$ и $b_1=b_0-\chi$.

(b) При растяжении группы $G$ в $f$ раз клеток становится больше на $(f-1)|G|$, склеек — на $(f-1)|G|+(f-1)r$, где $r$ — число склеек внутри $G$, перпендикулярных оси; каждая такая пара даёт лесенку $2\times f$ с $f-1$ целиком склеенными квадратами. Число кусков не меняется (П5). Изменение $d$ по (a) равно нулю.

(c) Склейка внутри куска увеличивает $|S|$ на 1 и $q$ на $0$, $1$ или $2$; склейка двух кусков уменьшает число кусков на 1 и не меняет $q$.

(d) Нижняя оценка следует из (b), (c) и $d=0$ для одной клетки. Для верхней растим объект вдоль остовного дерева графа склеек (каждая новая клетка — растяжение соседней, сразу склеенной с ней), после чего добавляем оставшиеся $|S|-|O|+1$ склеек; при $q=0$ это ровно $d$.
\end{proof}

Подсчёт по формуле (a) совпал с прямым вычислением гомологий на $10\,550$ случайных растяжениях и $337$ склейках. Следствие: у дыры два способа появиться — склейкой (ноль бит) или пробиванием (не меньше бита на клетку по теореме~\ref{thm:loss}).

\begin{remark}
В размерности 3 для объекта, где склеено всё касающееся, число сквозных дыр равно $b_0+P-\chi$, где $P$ — число полостей; это следует из формулы Эйлера—Пуанкаре и двойственности Александера и известно в цифровой топологии \cite{toriwaki,kaczynski}. Система приходит к этому счёту сама, но новым его не считаем.
\end{remark}

\section{Число действий}\label{sec:count}

Обозначим через $m(O)$ число осей, вдоль которых у клеток $O$ есть хотя бы два разных значения координаты.

\begin{theorem}\label{thm:bounds}
\begin{enumerate}[label=\upshape(\alph*)]
\item Для объекта $O$, связного по граням, $m(O)\le\ell(O)\le|O|-1$, причём верхняя оценка достигается одними растяжениями без потерь. Для прямоугольного бруска со сторонами $\ge2$ по $k$ осям $\ell=k$.
\item Для любого объекта $O$ при разрешённом $\sigma$ имеем $\ell(O)\le n+1$; соответствующая цепочка (брусок-рамка, затем вырезание $\sigma_{P\setminus O}$) теряет ровно $|P|-|O|$ бит, где $P$ — наименьший брусок, содержащий $O$.
\end{enumerate}
\end{theorem}
\begin{proof}
(a) Множество значений $i$-й координаты может вырасти только при $\alpha$ вдоль оси $i$: остальные действия его не увеличивают. Верхняя оценка — рост по одной клетке, как в доказательстве теоремы~\ref{thm:loss}. Брусок строится растяжением по каждой оси по разу. (b) Рамка — по (a); у $\sigma_Z$ ровно $2^{|Z|}$ прообразов.
\end{proof}

Пункт (b) выражает обмен «работы» на потерю: любая фигура строится за $n+1$ действий, но ценой бит за каждую вырезанную клетку.

\paragraph{Рельефы.} Пусть $B\subset\Z^{n-1}$ — основание, $h\colon B\to\{1,2,\dots\}$ — высоты и $O=\{(b,j):0\le j<h(b)\}$. Рассмотрим модель, в которой каждое действие растягивает вверх зону из целых столбиков с коэффициентом $f$; тогда высота столбика равна произведению коэффициентов действий, в зону которых он попадал. Пусть $H$ — множество различных высот больше 1, $D=|H|$, и $R(H)$ — наименьшее число множителей, из которых каждая высота собирается как произведение некоторой их части (каждый множитель — не более одного раза). Тогда $\ell(O)\le \ell(B)+R(H)\le\ell(B)+D$.

\begin{theorem}\label{thm:np}
Если все высоты свободны от квадратов, то $R(H)$ равно наименьшему размеру нормального базиса семейства множеств простых делителей высот. Следовательно, вычисление $R$ NP-трудно (при записи высот в двоичной системе).
\end{theorem}
\begin{proof}
Нормальный базис семейства $\mathcal F$ — такой набор множеств $\mathcal B$, что каждое множество из $\mathcal F$ есть объединение попарно непересекающихся множеств из $\mathcal B$. Множитель, делящийся на $p^2$, не входит в разложение бесквадратной высоты; остальным множителям $m$ сопоставим множества $P(m)$ их простых делителей. Если $h=m_1\cdots m_r$, то $P(m_j)$ попарно не пересекаются и дают в объединении $P(h)$. Обратно, по нормальному базису строим множители — произведения простых каждого множества. Размеры совпадают. Любое семейство подмножеств конечного множества кодируется так, если сопоставить элементам различные простые. Задача о нормальном базисе NP-полна в сильном смысле \cite{jiang}.
\end{proof}

\begin{theorem}\label{thm:R}
\begin{enumerate}[label=\upshape(\alph*)]
\item $\lceil\log_2(D+1)\rceil\le R(H)\le D$.
\item Если высоты попарно взаимно просты, то $R=D$.
\item Для $H=\{b,b^2,\dots,b^D\}$, $b\ge2$, имеем $R=\lceil\log_2(D+1)\rceil$.
\item Существует наименьший набор множителей, все элементы которого лежат в $C(H)=H\cup\mathcal D\cup\mathcal Q$, где $\mathcal D$ — числа $\ge2$, делящие хотя бы две высоты, а $\mathcal Q=\{h/d\ge2: h\in H,\ d\in\mathcal D,\ d\mid h\}$. Поэтому при фиксированном $k$ условие $R\le k$ проверяется за время, полиномиальное от $D$ и наибольшей высоты.
\end{enumerate}
\end{theorem}
\begin{proof}
(a) Из $R$ множителей получается не больше $2^R-1$ непустых частей, а разным высотам нужны разные части. (b) Множитель, общий для двух высот, делил бы обе. (c) Множители $b,b^2,b^4,\dots,b^{2^{k-1}}$ дают все $b^e$, $e\le 2^k-1$, по двоичной записи $e$. (d) В наименьшем наборе множитель, используемый двумя высотами, лежит в $\mathcal D$. У каждой высоты не больше одного «личного» множителя (два можно заменить произведением). Личный множитель высоты $h$ заменяем на $h/a$, если в разложении $h$ есть общий множитель $a$, и на $h$ иначе; размер набора не растёт. Перебор наборов из $k$ кандидатов, $|C(H)|=O(D^3h_{\max})$, полиномиален при фиксированном $k$.
\end{proof}

Таким образом, для числа действий существуют точные «лёгкие» случаи (b), (c) и «трудная середина»; для полного числа $\ell$ (с частичными зонами, где высоты не только умножаются, но и прибавляются) трудность не доказана.

\section{Гипотезы и открытые проблемы}\label{sec:conj}

\paragraph{Правило отрезков.} Для оси $i$ \emph{отрезком} объекта назовём максимальный ряд подряд идущих клеток вдоль оси $i$. Назовём \emph{шагом по отрезкам} действие одного из двух видов (ось $i$ и направление $s$ выбраны):
\begin{itemize}
\item \emph{умножить}: взять несколько отрезков из разных линий с одинаковым началом (крайней клеткой со стороны, обратной $s$) и растянуть их объединение в $f\ge2$ раз;
\item \emph{нарастить}: взять несколько отрезков из разных линий с одинаковым концом и растянуть их концевые клетки в $a+1$ раз ($a\ge1$), то есть нарастить каждый на $a$ клеток.
\end{itemize}
Каждый шаг — частный случай действия $\alpha$. Обозначим через $\ell_{\alpha}(O)$ наименьшее число растяжений с любыми зонами (при П3 и П5), а через $\ell_{\mathrm{seg}}(O)$ — наименьшее число шагов по отрезкам.

\begin{conjecture}[правило отрезков]\label{conj:seg}
Для любого объекта $O$, получаемого из одной клетки растяжениями, $\ell_{\mathrm{seg}}(O)=\ell_\alpha(O)$.
\end{conjecture}

Число отрезков линейно по $|O|$, тогда как зон $2^{|O|}$, поэтому гипотеза резко сокращает перебор (но по теореме~\ref{thm:np} не устраняет его). Проверка — полный перебор всех цепочек растяжений с любыми зонами против цепочек шагов по отрезкам:

\begin{center}
\begin{tabular}{lccr}
\toprule
Поле & Размерность & Длина цепочек & Фигур (все совпали)\\
\midrule
$3\times6$ & 2 & $\le3$ & $8\,631$\\
$4\times5$ & 2 & $\le3$ & $9\,031$\\
$3\times4$ & 2 & $\le4$ & $817$\\
$2\times3\times3$ & 3 & $\le3$ & $7\,196$\\
$3\times3\times3$ & 3 & $\le3$ & $30\,893$\\
$2\times2\times3$ & 3 & $\le4$ & $1\,670$\\
$2\times2\times4$ & 3 & $\le4$ & $16\,732$\\
$4\times4\times4$ & 3 & $\le3$ & $4\,760\,995$\\
\midrule
Всего & & & $4\,835\,965$\\
\bottomrule
\end{tabular}
\end{center}

Правило П5 существенно: если разрешить растяжению рвать склейки, на поле $3\times6$ достижимо $9870$ фигур против $8632$ (с учётом одной клетки), и правило отрезков их не описывает. На рельефах более простой вариант (шаги по граням) совпал с точным минимумом на всех $859$ проверенных рельефах.

\begin{conjecture}
Для рельефов на плоскости $\ell\ge1+\lceil\log_2 D'\rceil$, где $D'$ — число различных высот рельефа. (Проверено перебором на $1307$ рельефах.)
\end{conjecture}

\begin{conjecture}
В теореме~\ref{thm:holes}(d) ровно $d$ склеек достаточно и при $q>0$.
\end{conjecture}

\begin{problem}
Разрешима ли задача вычисления $\ell(O)$ в общем случае? Зоны и коэффициенты не ограничены, а промежуточные объекты могут быть сколь угодно большими.
\end{problem}

\begin{problem}
Для каких классов наборов высот $R(H)$ вычисляется за полиномиальное время? Существует ли приближённый алгоритм с постоянным множителем?
\end{problem}

\begin{problem}\label{pr:limit}
Существует ли предел цепочек, содержащих сжатия и $\sigma$, при измельчении сетки (ср. теорему~\ref{thm:affine}(c))?
\end{problem}

\begin{problem}
Какие утверждения о дырах (теорема~\ref{thm:holes}) переносятся на тонкие объекты, например на ленты из квадратов в $\Z^3$? (Лента Мёбиуса из квадратов решётки существует; наименьшая найденная состоит из 7 квадратов.)
\end{problem}

\section{Формальная проверка и программы}\label{sec:formal}

Ядро системы записано в Lean~4 \cite{lean} (версия 4.37.0-rc1, без сторонних библиотек): определения точек, объектов, действий $\alpha$, $\sigma$, $\beta$, правила П3 и цепочек и 36 теорем, опирающиеся только на стандартные аксиомы (\texttt{propext}, \texttt{Classical.choice}, \texttt{Quot.sound}), без пропусков (\texttt{sorry}). Машинно доказаны: нижняя оценка $m(O)\le\ell(O)$ теоремы~\ref{thm:bounds}(a) для всего ядра; отсутствие потерь при растяжении; прообразы $\sigma$; структурная часть леммы~\ref{lem:bit}; шаг роста теоремы~\ref{thm:loss}; неубывание рамки при растяжении (поэтому переборы в ограниченной коробке в §\ref{sec:conj} точны); ключевой шаг сведения теоремы~\ref{thm:np}; формула числа клеток при растяжении и отсутствие наложений; слияние клеток при сжатии; счёт вариантов в лемме~\ref{lem:bit}; пункты (a)–(c) теоремы~\ref{thm:R} и шаги замены из (d); обратимость поворота на $90^\circ$; клеточная часть теоремы~\ref{thm:affine}(a,b); подсчёт вершин по всем 256 конфигурациям и пункты (b)–(d) теоремы~\ref{thm:holes}. Не записаны машинно топологическая часть теоремы~\ref{thm:holes}(a), пределы и NP-трудность (последняя опирается на опубликованный результат \cite{jiang}).

Все переборы выполнены программами на Python и C++; программы, данные и файлы Lean открыты (ссылка в сноске на первой странице).

\section{Близкие работы}

\emph{Логическая необратимость.} Мера потери как $\log_2$ числа прообразов — стандартная \cite{landauer,frank}; новым является её применение к действиям над формами (лемма~\ref{lem:bit}, теоремы~\ref{thm:loss}, \ref{thm:exact}).
\emph{Математическая морфология} \cite{serra}: растяжение и сжатие близки к дилатации и эрозии, но там действуют структурным элементом на всё изображение, а не растягивают выбранную часть от её края; подсчёта потерь в битах там не найдено.
\emph{Аффинные преобразования}: теорема~\ref{thm:affine} показывает, что $\alpha$ — частный случай аффинного растяжения; различие — в способе задания и в учёте.
\emph{Цифровая топология и вычислительная гомология} \cite{kong,kaczynski,toriwaki}: формула теоремы~\ref{thm:holes}(a) — частный случай счёта эйлеровой характеристики; новое — что склеиваются не все касающиеся клетки.
\emph{Перестройка и самосборка фигур} \cite{akitaya,soloveichik}: ближайшие аналоги числа действий, но с иными элементарными операциями.
\emph{Теория сборки} \cite{assembly}: индекс сборки — наименьшее число соединений частей; число действий $\ell$ — аналог для геометрических растяжений, а потеря в битах в теории сборки не рассматривается.

\begin{thebibliography}{99}
\bibitem{landauer} R.~Landauer. Irreversibility and heat generation in the computing process. \emph{IBM J. Res. Dev.} \textbf{5}(3), 183--191, 1961.
\bibitem{frank} M.~P.~Frank. Physical foundations of Landauer's principle. arXiv:1901.10327, 2019.
\bibitem{serra} J.~Serra. \emph{Image Analysis and Mathematical Morphology}. Academic Press, 1982.
\bibitem{kong} T.~Y.~Kong, A.~Rosenfeld. Digital topology: introduction and survey. \emph{Comput. Vision Graph. Image Process.} \textbf{48}, 357--393, 1989.
\bibitem{kaczynski} T.~Kaczynski, K.~Mischaikow, M.~Mrozek. \emph{Computational Homology}. Springer, 2004.
\bibitem{toriwaki} J.~Toriwaki, H.~Yoshida. \emph{Fundamentals of Three-Dimensional Digital Image Processing}. Springer, 2009.
\bibitem{jiang} T.~Jiang, B.~Ravikumar. Minimal NFA problems are hard. \emph{SIAM J. Comput.} \textbf{22}(6), 1117--1141, 1993.
\bibitem{akitaya} H.~A.~Akitaya et al. Compacting squares: input-sensitive in-place reconfiguration of sliding squares. arXiv:2105.07997, 2021.
\bibitem{soloveichik} D.~Soloveichik, E.~Winfree. Complexity of self-assembled shapes. \emph{SIAM J. Comput.} \textbf{36}(6), 1544--1569, 2007.
\bibitem{assembly} A.~Sharma, D.~Czégel, M.~Lachmann, C.~P.~Kempes, S.~I.~Walker, L.~Cronin. Assembly theory explains and quantifies selection and evolution. \emph{Nature} \textbf{622}, 321--328, 2023.
\bibitem{lean} L.~de~Moura, S.~Ullrich. The Lean 4 theorem prover and programming language. In: \emph{CADE-28}, LNCS 12699, 625--635, 2021.
\end{thebibliography}

\end{document}
